% GENERATED FILE. Edit canonical lessons, then run npm run generate:books.
% canonical-source-hash: fnv1a64:16d890e40d8b8420
% canonical-lessons: TA-C129-cinkam, TA-C129-karati, TA-C129-kuranku, TA-C129-muyal, TA-C129-man2

\chapter{Lion, Bear, Monkey, Rabbit, Deer}
\label{ch:ta-lion-bear-monkey-rabbit-deer}
\hlchaptermodality{129}

\begin{chapteropening}
\textbf{By the end of this chapter:} \emph{I can say a lion, a bear, a monkey, a rabbit and a deer.}

Complete the last lesson of chapter 129: i can say a lion, a bear, a monkey, a rabbit and a deer.
\end{chapteropening}

\section[ciṅkam]{\ta{சிங்கம்} (ciṅkam) — a lion}

\label{lesson:TA-C129-cinkam}



\begin{quote}
Before the new one: say the Tamil for an elephant, then the Tamil for a tiger.
\end{quote}

\subsection*{You'll want to know: \ta{சிங்கம்}}
\textbf{\ta{சிங்கம்}} — \emph{ciṅkam} — \textquotedblleft{}a lion\textquotedblright{}.

\begin{grammarlens}[title={What you've built}]
One more word for this chapter.
\end{grammarlens}

\subsection*{Your turn}
\begin{itemize}
\raggedright
  \item \emph{Say it:} \emph{ciṅkam}
  \item \emph{Say it:} \emph{ciṅkam}, once more
  \item \emph{Say it:} say \emph{ciṅkam}
  \item \emph{Recall:} say the Tamil for an elephant, then the Tamil for a tiger, then say \emph{ciṅkam} again
\end{itemize}

\begin{tcolorbox}[breakable,colback=teal!4,colframe=teal!35!black,title={Before you move on}]
What does \ta{சிங்கம்} mean? (\textquotedblleft{}A lion\textquotedblright{}.) Say it once more.
\end{tcolorbox}
\section[karaṭi]{\ta{கரடி} (karaṭi) — a bear}

\label{lesson:TA-C129-karati}



\begin{quote}
Before the new one: say the Tamil for a tiger, then the Tamil for a lion.
\end{quote}

\subsection*{You'll want to know: \ta{கரடி}}
\textbf{\ta{கரடி}} — \emph{karaṭi} — \textquotedblleft{}a bear\textquotedblright{}.

\begin{grammarlens}[title={What you've built}]
One more word for this chapter.
\end{grammarlens}

\subsection*{Your turn}
\begin{itemize}
\raggedright
  \item \emph{Say it:} \emph{karaṭi}
  \item \emph{Say it:} \emph{karaṭi}, once more
  \item \emph{Say it:} say \emph{karaṭi}
  \item \emph{Recall:} say the Tamil for a tiger, then the Tamil for a lion, then say \emph{karaṭi} again
\end{itemize}

\begin{tcolorbox}[breakable,colback=teal!4,colframe=teal!35!black,title={Before you move on}]
What does \ta{கரடி} mean? (\textquotedblleft{}A bear\textquotedblright{}.) Say it once more.
\end{tcolorbox}
\section[kuraṅku]{\ta{குரங்கு} (kuraṅku) — a monkey}

\label{lesson:TA-C129-kuranku}



\begin{quote}
Before the new one: say the Tamil for a lion, then the Tamil for a bear.
\end{quote}

\subsection*{You'll want to know: \ta{குரங்கு}}
\textbf{\ta{குரங்கு}} — \emph{kuraṅku} — \textquotedblleft{}a monkey\textquotedblright{}.

\begin{grammarlens}[title={What you've built}]
One more word for this chapter.
\end{grammarlens}

\subsection*{Your turn}
\begin{itemize}
\raggedright
  \item \emph{Say it:} \emph{kuraṅku}
  \item \emph{Say it:} \emph{kuraṅku}, once more
  \item \emph{Say it:} say \emph{kuraṅku}
  \item \emph{Recall:} say the Tamil for a lion, then the Tamil for a bear, then say \emph{kuraṅku} again
\end{itemize}

\begin{tcolorbox}[breakable,colback=teal!4,colframe=teal!35!black,title={Before you move on}]
What does \ta{குரங்கு} mean? (\textquotedblleft{}A monkey\textquotedblright{}.) Say it once more.
\end{tcolorbox}
\section[muyal]{\ta{முயல்} (muyal) — a rabbit}

\label{lesson:TA-C129-muyal}



\begin{quote}
Before the new one: say the Tamil for a bear, then the Tamil for a monkey.
\end{quote}

\subsection*{You'll want to know: \ta{முயல்}}
\textbf{\ta{முயல்}} — \emph{muyal} — \textquotedblleft{}a rabbit\textquotedblright{}.

\begin{grammarlens}[title={What you've built}]
One more word for this chapter.
\end{grammarlens}

\subsection*{Your turn}
\begin{itemize}
\raggedright
  \item \emph{Say it:} \emph{muyal}
  \item \emph{Say it:} \emph{muyal}, once more
  \item \emph{Say it:} say \emph{muyal}
  \item \emph{Recall:} say the Tamil for a bear, then the Tamil for a monkey, then say \emph{muyal} again
\end{itemize}

\begin{tcolorbox}[breakable,colback=teal!4,colframe=teal!35!black,title={Before you move on}]
What does \ta{முயல்} mean? (\textquotedblleft{}A rabbit\textquotedblright{}.) Say it once more.
\end{tcolorbox}
\section[māṉ]{\ta{மான்} (māṉ) — a deer}

\label{lesson:TA-C129-man2}



\begin{quote}
Before the new one: say the Tamil for a monkey, then the Tamil for a rabbit.
\end{quote}

\subsection*{You'll want to know: \ta{மான்}}
\textbf{\ta{மான்}} — \emph{māṉ} — \textquotedblleft{}a deer\textquotedblright{}.

\begin{grammarlens}[title={What you've built}]
One more word for this chapter.
\end{grammarlens}

\subsection*{Your turn}
\begin{itemize}
\raggedright
  \item \emph{Say it:} \emph{māṉ}
  \item \emph{Say it:} \emph{māṉ}, once more
  \item \emph{Say it:} say \emph{māṉ}
  \item \emph{Recall:} say the Tamil for a monkey, then the Tamil for a rabbit, then say \emph{māṉ} again
\end{itemize}

\begin{tcolorbox}[breakable,colback=teal!4,colframe=teal!35!black,title={Before you move on}]
What does \ta{மான்} mean? (\textquotedblleft{}A deer\textquotedblright{}.) Say it once more.
\end{tcolorbox}
